%-----------------------------------------------------------------------
\chapter{Implementação da Solução e Resultados}

\section{Visão geral da implementação}

O protótipo implementado consiste em uma aplicação web para gerenciamento de inventário patrimonial com apoio de eventos RFID. A solução permite cadastrar locais, leitores RFID e itens patrimoniais, associando cada item a uma etiqueta. Além disso, disponibiliza recursos para acionar janelas de leitura, registrar eventos, atualizar a localização física de ativos, identificar inconsistências e consultar o histórico operacional.

A implementação foi organizada de forma modular. O frontend concentra as telas de interação com o usuário, incluindo painel inicial, configurações, patrimônio, leitores RFID, auditoria, inconsistências e log operacional. O backend disponibiliza uma API REST responsável por persistência, validação e aplicação das regras de negócio. Na validação com o leitor de proximidade, utilizou-se um adaptador intermediário de integração capaz de capturar a tag lida pelo dispositivo e encaminhar o evento ao sistema.

Esse adaptador foi necessário porque o leitor empregado opera como dispositivo de entrada local e não possui comunicação HTTP nativa. Assim, sua utilização decorre de uma limitação de interface do equipamento disponível para o experimento, e não de uma restrição da API proposta. Essa organização permite que o leitor utilizado na validação seja substituído futuramente por outros dispositivos de captura. A condição para essa substituição é que o novo dispositivo, ou um gateway associado a ele, envie à API REST eventos compatíveis com o formato esperado. Portanto, embora a validação tenha sido realizada com leitor RFID de proximidade de baixo custo, a arquitetura não fica restrita a esse equipamento.

\section{Cadastro de locais, leitores e itens}

O cadastro de locais representa os espaços nos quais os bens podem estar logicamente alocados ou fisicamente detectados. Cada local possui identificação própria e é utilizado como referência para leitores RFID e itens patrimoniais. Essa modelagem permite comparar a localização esperada de um item com a localização detectada em uma leitura.

Os leitores RFID são cadastrados com identificador de hardware, local associado e tipo de operação. A distinção entre leitores do tipo destino e leitores do tipo fluxo permite representar comportamentos distintos. Um leitor de destino pode atualizar a localização física de um item, enquanto um leitor de fluxo pode apenas registrar a passagem de uma tag por determinado ponto. Essa decisão de modelagem é relevante para uma arquitetura escalável, pois possibilita que diferentes pontos de captura tenham funções distintas no processo de inventário.

Os itens patrimoniais são cadastrados com tag RFID, nome, local lógico, local físico e estado de atividade. A tag funciona como vínculo entre o objeto físico e sua representação no sistema. Quando uma leitura é recebida, o backend verifica se a tag está cadastrada e decide qual ação deve ser executada, como atualização de local, geração de histórico ou registro de inconsistência.

\section{Fluxo de eventos RFID}

O fluxo principal inicia-se com a ativação de uma janela de leitura. Essa janela representa o intervalo no qual o sistema espera receber eventos de determinado leitor. Em uma implantação escalável, essa ativação poderia ser causada por um sensor infravermelho ou outro sensor de presença, que acionaria a API para habilitar uma antena RFID em determinado ponto. No protótipo validado, a leitura foi realizada por um dispositivo de proximidade, e o adaptador intermediário de integração ficou responsável por converter a tag capturada em requisição HTTP enviada ao backend.

Ao receber um evento RFID, a API executa etapas de validação e processamento. Inicialmente, verifica a autenticidade da requisição e a existência do leitor. Em seguida, normaliza as tags recebidas, trata leituras duplicadas em curto intervalo e verifica se a leitura ocorreu durante uma janela ativa. Quando a tag corresponde a um item cadastrado, o sistema registra a leitura e atualiza o estado do item conforme o tipo do leitor. Quando a tag não está cadastrada, o sistema pode registrar uma inconsistência de tag desconhecida.

Esse fluxo evidencia a separação entre hardware e regra de negócio. O leitor é responsável apenas por capturar a identificação; a API concentra a interpretação do evento. Essa separação é um dos elementos que tornam a aplicação extensível, pois diferentes dispositivos podem atuar como fonte de eventos sem exigir reescrita completa das regras do sistema.

\section{Auditoria, inconsistências e log operacional}

A funcionalidade de auditoria permite comparar os itens esperados em um local com as tags efetivamente detectadas durante uma janela de leitura. A partir dessa comparação, o sistema pode classificar situações como item encontrado, item ausente, item em local divergente ou tag desconhecida. Essa classificação auxilia o processo de conferência patrimonial, pois transforma leituras isoladas em informações operacionais.

As inconsistências geradas pelo sistema registram divergências entre o inventário lógico e o inventário físico. Um item pode estar logicamente associado a determinado local, mas ser detectado em outro. Também pode ocorrer de uma tag ser lida sem estar cadastrada, indicando necessidade de investigação. O registro dessas inconsistências permite que o usuário acompanhe pendências e resolva divergências de forma organizada.

O log operacional mantém o histórico dos eventos relevantes do sistema. Esse histórico inclui movimentações, leituras, auditorias e alterações de estado. Em um contexto de controle patrimonial, esse recurso é importante porque permite rastrear ações e compreender como o estado de um item foi modificado ao longo do tempo.

\section{Interface do sistema}

A interface web foi organizada para apoiar o fluxo de uso do sistema. O painel inicial apresenta uma visão geral do inventário e das atividades recentes. A área de configurações permite cadastrar locais, leitores e itens patrimoniais. A tela de leitores permite acionar janelas de sincronização e auditoria. A tela de patrimônio apresenta os itens cadastrados e suas localizações. A tela de inconsistências permite consultar divergências detectadas, enquanto o log operacional centraliza o histórico de eventos.

Para documentação do TCC, recomenda-se utilizar capturas de tela do dashboard, da tela de configurações, da listagem de leitores, da auditoria, da tela de inconsistências, da consulta de patrimônio e do log operacional. Essas evidências ajudam a demonstrar que a solução não se limita ao processamento interno da API, mas oferece interface de uso para acompanhamento do inventário.

\section{Resultados da validação funcional}

A validação funcional foi realizada em ambiente controlado, utilizando um computador pessoal, um leitor RFID de proximidade de baixo custo e uma única tag RFID. O leitor empregado exige aproximação direta da tag para captura do identificador, diferindo de uma infraestrutura RFID de maior alcance. Assim, os testes tiveram como objetivo verificar o fluxo lógico da aplicação: leitura da tag, captura pelo adaptador intermediário, envio do evento à API REST, processamento pelo backend, persistência dos dados e atualização das informações de inventário. Não foram avaliados alcance, leitura simultânea de múltiplas tags ou desempenho físico de antenas RFID. A Tabela \ref{tab:testes-funcionais} resume os principais cenários de validação previstos para o sistema, organizados a partir da relação entre procedimento executado, resultado esperado e resultado observado, conforme práticas de documentação de testes de software \cite{iso29119}.

\begin{table}[H]
\centering
\caption{Cenários de validação funcional do protótipo}
\label{tab:testes-funcionais}
\begin{tabular}{|p{3.2cm}|p{4.1cm}|p{4.1cm}|p{2.7cm}|}
\hline
\textbf{Cenário} & \textbf{Procedimento} & \textbf{Resultado esperado} & \textbf{Resultado observado} \\
\hline
Cadastro básico & Cadastrar local, leitor RFID e item patrimonial com a tag disponível no experimento. & Registros disponíveis para consulta e uso nos fluxos de leitura. & A preencher com evidência do teste. \\
\hline
Leitura válida & Ativar janela de leitura e aproximar a tag cadastrada do leitor de proximidade. & Sistema registra a leitura e atualiza o local físico conforme o leitor. & A preencher com evidência do teste. \\
\hline
Tag desconhecida & Enviar código de tag não cadastrado, caso o adaptador ou a API sejam utilizados para simulação controlada. & Sistema registra evento e sinaliza inconsistência de tag desconhecida. & Cenário dependente de simulação ou de segunda tag; preencher somente se executado. \\
\hline
Local divergente & Associar o item a local lógico diferente do local vinculado ao leitor e aproximar a tag cadastrada. & Sistema identifica divergência entre localização lógica e física. & A preencher com evidência do teste. \\
\hline
Duplicidade & Aproximar a mesma tag repetidamente do leitor em curto intervalo. & Sistema ignora ou consolida leituras duplicadas conforme regra de deduplicação. & A preencher com evidência do teste. \\
\hline
Auditoria & Executar janela de auditoria para um local utilizando a tag disponível no experimento. & Sistema classifica o item lido e permite identificar ausências conforme os itens cadastrados para o local. & Validação parcial com uma tag; preencher com evidência do teste executado. \\
\hline
Leitor offline & Interromper o adaptador intermediário, desconectar o leitor ou simular falha de envio de evento. & Sistema bloqueia, sinaliza ou registra falha de comunicação sem comprometer os dados. & A preencher com evidência do teste. \\
\hline
\end{tabular}
\end{table}

Como os resultados quantitativos dependem de medições específicas do experimento, não foram assumidos valores de acurácia, alcance, taxa de leitura ou tempo de resposta sem evidência. A validação apresentada deve ser complementada com registros reais, como prints, logs e tabelas de execução. Ainda assim, a implementação demonstra que o fluxo essencial da solução foi contemplado: o sistema cadastra entidades, recebe eventos RFID, processa leituras, atualiza o inventário, gera histórico e identifica inconsistências. Essa forma de avaliação deve ser entendida como validação funcional do protótipo no escopo definido, e não como certificação de desempenho físico da tecnologia RFID em ambiente real de larga escala \cite{ieee1012,knapp2022}.

Portanto, o resultado principal do protótipo é a comprovação funcional da arquitetura de software. A validação com computador pessoal, leitor de proximidade e uma tag demonstra que a aplicação é capaz de operar com eventos RFID reais em ambiente controlado, comprovando o fluxo computacional necessário ao inventário com RFID. A evolução para sensores infravermelhos e antenas RFID de maior alcance permanece como continuidade natural da proposta, exigindo novos testes físicos de alcance, taxa de leitura, interferência e leitura simultânea.
